\begin{table}[htbp]
\centering
\caption{Comparação PNADC vs.\ INEP -- Transição 2019/2020, Brasil}
\label{tab:t4_pnadc_inep}
\begin{threeparttable}
\small
\begin{tabular}{lccccc}
\toprule
                       & \multicolumn{3}{c}{Indicador (\%)} & \multicolumn{2}{c}{Gap} \\
\cmidrule(lr){2-4} \cmidrule(lr){5-6}
                       & Promoção & Repetência & Evasão & PP & Relativo \\
\midrule
\multicolumn{6}{l}{\textit{EF iniciais (Anos Iniciais)}} \\
INEP                   & 93,5 & 5,1 & 1,1  & --   & --    \\
PNADC                  & 76,6 & 22,9 & 0,3 & $-$16,9 & $-$18\% \\
\midrule
\multicolumn{6}{l}{\textit{EF finais (Anos Finais)}} \\
INEP                   & 86,0 & 9,8 & 3,5  & --   & --    \\
PNADC                  & 71,5 & 24,7 & 1,6 & $-$14,5 & $-$17\% \\
\midrule
\multicolumn{6}{l}{\textit{Ensino Médio (Total)}} \\
INEP                   & 82,7 & 8,3 & 6,9  & --   & --    \\
PNADC                  & 49,0 & 17,3 & 24,5 & $-$33,7 & $-$41\% \\
\bottomrule
\end{tabular}
\begin{tablenotes}
\footnotesize
\item \textit{Nota:} Fontes -- INEP Taxa de Transição 2019/2020 (Brasil),
publicada em \citet{INEP_transicao}; PNADC trimestral, painel longitudinal,
transição 2019$\to$2020 medida via primeira observação em $t$ vs.\ primeira
observação em $t+1$. Gap em pontos percentuais (PP) e em \% do nível INEP.
A diferença substancial para o EM reflete principalmente a definição
operacional de promoção: o INEP classifica como promovidos os alunos que
concluíram o 3\textsuperscript{o} EM mesmo sem matrícula subsequente;
nossa metodologia inicial exige frequência escolar em $t+1$. Versão refinada
da PNADC em desenvolvimento.
\end{tablenotes}
\end{threeparttable}
\end{table}
